%%
%% (
%%  )\ )                             (
%%  (()/(   (            (             )\  )   (
%%   /(_))  ))\   (       ))\  (   (   (()/(   ))\
%%   (_))  /((_)  )\  )  /((_) )\  )\   ((_))/((_)
%%   | _ \(_))(  _(_/( (_) )  ((_)((_)  _| |(_))
%%   |   /| || || ' \))/ -_)/ _|/ _ \/ _` |/ -_)
%%   |_|_\ \_,_||_||_| \___|\__|\___/\__,_|\___|
%%

\documentclass{article}
\usepackage[utf8]{inputenc}
\usepackage{amsmath}
\usepackage{amsfonts}
\usepackage{amssymb}
\usepackage{graphicx} % Paquete para incluir imágenes en el documento LaTeX
\usepackage{hyperref}
\hypersetup{
  colorlinks=true,
  linkcolor=blue,
  filecolor=magenta,
  urlcolor=cyan,
}
\urlstyle{same}
\usepackage{varwidth}

\newcommand\tab[1][1cm]{\hspace*{#1}}

\usepackage{multirow}

\usepackage[a4paper,rmargin=1.5cm,lmargin=1.5cm,top=1.5cm,bottom=1.5cm]{geometry}

\usepackage{pdfpages}

\usepackage{xcolor}
\usepackage{minted}
\usemintedstyle[cpp]{monokai}
\usemintedstyle[python]{paraiso-dark}
\definecolor{LightGray}{gray}{0.9}
\definecolor{DarkGray}{gray}{0.1}
\definecolor{MidGray}{gray}{0.8}
\definecolor{codegreen}{rgb}{0,0.6,0}
\definecolor{codegray}{rgb}{0.5,0.5,0.5}
\definecolor{codepurple}{rgb}{0.58,0,0.82}
\definecolor{backcolour}{rgb}{0.95,0.95,0.92}

\setlength{\parindent}{0px}  % Setea la indentacion de la primera linea de cada parrafo a cero pixeles.


\title{Resolución de la quinta semana}
\author{@RuneCode}

\begin{document}
%% Portada
\includepdf{./portada/portada.pdf}


%-----------------------------------------------------------------------------------------
%Primer problema
\section{Quinta Semana}%
En esta clase se revisaron los temas de:
\begin{itemize}
\item \textbf{Uso de apuntadores}
\item \textbf{Programación Modular}
\item \textbf{Tipos de Subprogramas}(Función, Predicado y Procedimiento)
\end{itemize}
\vspace{1cm}
\textbf{Nota:} Te recomendamos realizar los ejercicios antes de ver su solución
ya que el objetivo de este documento es presentarte una de las muchas posibles
soluciones. Si tienes alguna sugerencia de mejora puedes comunicarte por correo
a la dirección \href{mailto:gprunecode@gmail.com}{gprunecode@gmail.com}.

\section{Ejemplos}%
\textbf{\textit{Problema 1}} Implementar una \textbf{función} para la suma de 2 números enteros\\
\textit{Resolución en pseudocódigo}\\ 

%Pseudocódigo
\inputminted
[
  frame=lines,
  framesep=2mm,
  baselinestretch=1.2,
  rulecolor=\color{black!30},
  %fontsize=\footnotesize,
  bgcolor=LightGray,
  %linenos
]
{./pseudocode.py:PseudocodeLexer -x}
{./pseudocodigo/001.algo}

%C
\textit{Resolución en C}\\
\inputminted
[
  frame=lines,
  framesep=2mm,
  baselinestretch=1.2,
  rulecolor=\color{black!30},
  %fontsize=\footnotesize,
  bgcolor=DarkGray,
  %linenos
]
{cpp}
{./c/ejercicio1.cpp}
\clearpage

\vspace{1cm}

%---------------------------------------------------------------------------------
%Segundo problema

\textbf{\textit{Problema 2}} Implementar un \textbf{procedimiento} para la suma de 2 números enteros\\
\textit{Resolución en pseudocódigo}\\ 

%Pseudocódigo
\inputminted
[
  frame=lines,
  framesep=2mm,
  baselinestretch=1.2,
  rulecolor=\color{black!30},
  %fontsize=\footnotesize,
  bgcolor=LightGray,
  %linenos
]
{./pseudocode.py:PseudocodeLexer -x}
{./pseudocodigo/002.algo}


%C
\textit{Resolución en C}\\

\inputminted
[
  frame=lines,
  framesep=2mm,
  baselinestretch=1.2,
  rulecolor=\color{black!30},
  %fontsize=\footnotesize,
  bgcolor=DarkGray,
  %linenos
]
{cpp}
{./c/ejercicio2.cpp}
\clearpage

\vspace{1cm}

%------------------------------------------------------------------------------------
%Tercer problema

\textbf{\textit{Problema 3}} Hallar la suma de los n primeros multiplos de 3\\
\textit{Resolución en pseudocódigo}\\ 

%Pseudocódigo
\inputminted
[
  frame=lines,
  framesep=2mm,
  baselinestretch=1.2,
  rulecolor=\color{black!30},
  %fontsize=\footnotesize,
  bgcolor=LightGray,
  %linenos
]
{./pseudocode.py:PseudocodeLexer -x}
{./pseudocodigo/003.algo}


%C
\textit{Resolución en C}\\

\inputminted
[
  frame=lines,
  framesep=2mm,
  baselinestretch=1.2,
  rulecolor=\color{black!30},
  %fontsize=\footnotesize,
  bgcolor=DarkGray,
  %linenos
]
{cpp}
{./c/ejercicio3.cpp}
\clearpage

\vspace{1cm}

%---------------------------------------------------------------------------------
%Cuarto problema

\textbf{\textit{Problema 4}} Implemente un \textbf{procedimiento} que intercambie los valores de a y b\\

\textit{Resolución en pseudocódigo}\\ 

%Pseudocódigo
\inputminted
[
  frame=lines,
  framesep=2mm,
  baselinestretch=1.2,
  rulecolor=\color{black!30},
  %fontsize=\footnotesize,
  bgcolor=LightGray,
  %linenos
]
{./pseudocode.py:PseudocodeLexer -x}
{./pseudocodigo/004.algo}


%C
\textit{Resolución en C}\\
\inputminted
[
  frame=lines,
  framesep=2mm,
  baselinestretch=1.2,
  rulecolor=\color{black!30},
  %fontsize=\footnotesize,
  bgcolor=DarkGray,
  %linenos
]
{cpp}
{./c/ejercicio4.cpp}
\clearpage

\vspace{1cm}

%---------------------------------------------------------------------------------
\vspace{3cm} 
\section{¡Envianos tus soluciones!}
Si estás llevando este curso con los profesores Cabrera, Romero o Salinas; envíanos tus soluciones en los diferentes lenguajes de programación que conozcas al correo \href{mailto:gprunecode@gmail.com}{gprunecode@gmail.com}. n.n \\ 


Las mejores soluciones tanto en el algoritmo como en el código, serán publicadas en las siguientes ediciones de estos documentos.\\ \\

El asunto del correo debe estar de la siguiente manera:\\
$NumeroDeSesion-Profesor-NumeroDeEjercicio$ \\ \\
Por ejemplo:  \\
$01-Romero-01$ \\

Y dentro del correo adjuntar tu solución y nombre como quieras ser reconocido en caso de ser electo.

\vspace{2cm}
\LARGE\textit{RuneCode}

\end{document}
